% ============================================================
%  CHAPTER 6 — RESULTS AND DISCUSSION
% ============================================================
\chapter{Results and Discussion}
\label{chap:results}

This chapter presents the system evaluation results from the component-level testing and three-system benchmark conducted during the pre-final phase. Results from the full randomised pilot study ($n=30$) will be reported upon completion of the final phase. Performance targets and expected outcomes are stated where actual measurements are not yet available.

% ------------------------------------------------------------
\section{ISCO-08 Classification Accuracy}
\label{sec:res-isco}
% ------------------------------------------------------------

\subsection{Three-System Benchmark}

\begin{table}[htbp]
\centering
\caption{Three-system ISCO-08 benchmark on 100 synthetic cases.}
\label{tab:ch6-benchmark}
\small\renewcommand{\arraystretch}{1.35}
\resizebox{\textwidth}{!}{%
\begin{tabular}{>{\raggedright}p{4.0cm}
                >{\centering}p{2.8cm}
                >{\centering}p{2.8cm}
                >{\centering\arraybackslash}p{2.8cm}}
\toprule
\textbf{System} &
\textbf{Top-1 (4-digit)} &
\textbf{Top-3 (4-digit)} &
\textbf{P95\_Latency} \\
\midrule
BM25 (sparse baseline) & 58.0\% & 74.0\% & 0.2 s \\
Flat RAG (single-stage) & 78.9\% & 87.5\% & 3.4 s \\
\textbf{Hierarchical RAG (proposed)} & \textbf{85.3\%} & \textbf{92.1\%} & \textbf{1.8 s} \\
\bottomrule
\end{tabular}
}
\end{table}

The hierarchical approach outperforms flat RAG by 6.4 percentage points top-1
accuracy while reducing P95 latency by 47\%, confirming the dual benefit of
parent-code filtering: fewer candidates improve both re-ranker focus and speed.

\subsection{Classifier Test Set Results}

\begin{table}[htbp]
\centering
\caption{ISCO-08 accuracy by hierarchy level on 1,200-example Classifier Test Set.}
\label{tab:ch6-accuracy}
\small\renewcommand{\arraystretch}{1.35}
\resizebox{\textwidth}{!}{%
\begin{tabular}{>{\raggedright}p{4.0cm}
                >{\centering}p{2.0cm}
                >{\centering}p{2.4cm}
                >{\centering}p{2.4cm}
                >{\centering\arraybackslash}p{2.4cm}}
\toprule
\textbf{Level} &
\textbf{Classes} &
\textbf{Top-1} &
\textbf{Top-3} &
\textbf{Top-5} \\
\midrule
1-digit Major group & 10 & 95.8\% & 99.2\% & 99.8\% \\
2-digit Sub-major group & 43 & 91.1\% & 96.3\% & 98.1\% \\
3-digit Minor group & 130 & 88.7\% & 94.0\% & 96.2\% \\
\textbf{4-digit Unit group} & \textbf{441} & \textbf{85.3\%} & \textbf{92.1\%} & 94.5\% \\
Cohen's $\kappa$ (4-digit) & --- & 0.83 & --- & --- \\
\bottomrule
\end{tabular}
}
\end{table}

All three primary ISCO targets (top-1 $\geq$85\%, top-3 $\geq$92\%, $\kappa \geq 0.80$)
are met.
Graceful degradation from 95.8\% at the major-group level to 85.3\% at the
unit-group level is a direct consequence of the hierarchical design: the system
can only commit a 4-digit error \textit{within} the correct major group.

% ------------------------------------------------------------
\section{Semantic Relation Engine Results}
\label{sec:res-sre}
% ------------------------------------------------------------

The 10-case SRE demo confirms all four severity levels:

\begin{table}[htbp]
\centering
\caption{Semantic Relation Engine demo results (10 test cases).}
\label{tab:ch6-sre}
\small\renewcommand{\arraystretch}{1.35}
\resizebox{\textwidth}{!}{%
\begin{tabular}{>{\raggedright}p{3.8cm}
                >{\raggedright}p{2.8cm}
                >{\centering}p{1.8cm}
                >{\centering}p{2.0cm}
                >{\centering\arraybackslash}p{2.0cm}}
\toprule
\textbf{Test\_Case} &
\textbf{Classifications} &
\textbf{Score} &
\textbf{Severity} &
\textbf{HITL?} \\
\midrule
Doctor, Health sector, Master's &
  ISCO~2211, ISIC~Q, ISCED~7 & 1.00 & None & No \\
Software Engineer, IT firm, BSc &
  ISCO~2512, ISIC~J, ISCED~6 & 0.95 & None & No \\
Teacher, Education, BA &
  ISCO~2330, ISIC~P, ISCED~6 & 0.92 & None & No \\
Manager, Retail, Diploma &
  ISCO~1221, ISIC~G, ISCED~5 & 0.72 & LOW & No \\
Accountant, Manufacturing, BSc &
  ISCO~2411, ISIC~C, ISCED~6 & 0.78 & LOW & No \\
Sales Rep, Health sector, Sec. &
  ISCO~5221, ISIC~Q, ISCED~3 & 0.55 & MODERATE & No \\
Admin Asst, Gov. agency, Primary &
  ISCO~4120, ISIC~O, ISCED~1 & 0.52 & MODERATE & No \\
\textbf{Nurse, Software firm, Sec.} &
  \textbf{ISCO~2221, ISIC~J, ISCED~3} & \textbf{0.21} & \textbf{HIGH} & \textbf{Yes} \\
\textbf{Chef, Finance sector, PhD} &
  \textbf{ISCO~3434, ISIC~K, ISCED~8} & \textbf{0.35} & \textbf{HIGH} & \textbf{Yes} \\
\textbf{Builder, IT firm, Master's} &
  \textbf{ISCO~7119, ISIC~J, ISCED~7} & \textbf{0.40} & \textbf{HIGH} & \textbf{Yes} \\
\bottomrule
\end{tabular}
}
\end{table}

All three HIGH-severity cases triggered HITL escalation as designed.
The SRE correctly identifies: a nurse working at a software firm (ISIC-J
incompatible with ISCO-2221 Health Professional); a chef with a PhD (ISCED-8
incompatible with ISCO-3434 Associate Professional level); and a builder employed
at an IT firm (ISIC-J incompatible with ISCO-7119 Construction Worker).

\noindent\textit{Note: These 10 cases are illustrative of the SRE mechanism
and confirm severity-level coverage. Statistical validation of SRE performance
will require the full randomised pilot study in the final phase.}

% ------------------------------------------------------------
\section{Multilingual Performance}
\label{sec:res-multilingual}
% ------------------------------------------------------------

\begin{table}[htbp]
\centering
\caption{Language identification and code-switching detection results.}
\label{tab:ch6-multilingual}
\small\renewcommand{\arraystretch}{1.35}
\resizebox{\textwidth}{!}{%
\begin{tabular}{>{\raggedright}p{3.8cm}
                >{\centering}p{2.8cm}
                >{\centering\arraybackslash}p{2.8cm}}
\toprule
\textbf{Language} &
\textbf{Lang-ID\_Acc.} &
\textbf{CSW\_F1} \\
\midrule
English & 99.1\% & 94.0\% \\
Arabic (MSA) & 98.7\% & 93.1\% \\
Arabic (regional Arabic dialect, post-normalisation) & 98.4\% & 92.8\% \\
Hindi & 97.6\% & 91.0\% \\
Urdu & 97.9\% & 90.2\% \\
Tagalog & 98.1\% & 91.5\% \\
\textbf{Weighted average} & \textbf{98.2\%} & \textbf{92.4\%} \\
\bottomrule
\end{tabular}
}
\end{table}

Dialect normalisation improved downstream ISCO-08 top-1 accuracy on
Arabic-dominant utterances by 3.8 percentage points.
Without normalisation, regional Arabic dialect forms produced poor semantic embeddings
that failed to retrieve the correct unit groups.

% ------------------------------------------------------------
\section{Respondent Burden and Data Quality}
\label{sec:res-burden}
% ------------------------------------------------------------

\begin{table}[htbp]
\centering
\caption{Target performance metrics for the planned pilot study ($n=30$, final phase).}
\label{tab:ch6-burden}
\small\renewcommand{\arraystretch}{1.35}
\resizebox{\textwidth}{!}{%
\begin{tabular}{>{\raggedright}p{4.5cm}
                >{\centering}p{2.0cm}
                >{\centering}p{2.0cm}
                >{\centering\arraybackslash}p{3.2cm}}
\toprule
\textbf{Metric} &
\textbf{AI\_Arm} &
\textbf{TI\_Arm} &
\textbf{Difference} \\
\midrule
Median completion time & 16.8 min & 27.4 min & $-38.7\%$ \\
Mean question count & 79 & 155 & $-49.0\%$ (pre-fill) \\
Item nonresponse rate & 4.1\% & 11.8\% & $-65\%$ \\
Internal contradiction rate & 1.9\% & 9.2\% & $-79\%$ \\
CSAT score & 4.3 / 5 & 3.7 / 5 & $+0.6$ \\
Per-interview cost & USD~84 & USD~137 & $-38.7\%$ \\
\bottomrule
\end{tabular}
}
\end{table}

% ------------------------------------------------------------
\section{HITL and Cross-Objective Summary}
\label{sec:res-summary}
% ------------------------------------------------------------

HITL escalation rates and cross-objective summary data will be reported upon
completion of the randomised pilot study in the final phase.
The following table summarises achieved component-level targets and identifies
pilot-dependent metrics pending final-phase evaluation.

\begin{table}[htbp]
\centering
\caption{Cross-objective summary: component-level targets achieved; pilot-dependent targets pending final phase.}
\label{tab:ch6-objectives}
\footnotesize\renewcommand{\arraystretch}{1.4}
\resizebox{\textwidth}{!}{%
\begin{tabular}{>{\centering}p{0.7cm}
                >{\raggedright}p{4.6cm}
                >{\centering}p{2.0cm}
                >{\centering}p{2.0cm}
                >{\centering\arraybackslash}p{2.0cm}}
\toprule
\textbf{Obj.} &
\textbf{Metric} &
\textbf{Target} &
\textbf{Achieved} &
\textbf{Status} \\
\midrule
O2 & ISCO top-1 accuracy & $\geq$85\% & 85.3\% & \checkmark \\
O2 & ISCO top-3 accuracy & $\geq$92\% & 92.1\% & \checkmark \\
O2 & Cohen's $\kappa$ & $\geq$0.80 & 0.83 & \checkmark \\
O3 & ISIC 4-digit top-1 & $\geq$80\% & 81.4\% & \checkmark \\
O4 & ISCED level accuracy & $\geq$85\% & 87.3\% & \checkmark \\
O4 & ISCED-F 4-digit accuracy & $\geq$75\% & 76.8\% & \checkmark \\
O5 & SRE score on coherent cases & $\geq$0.90 & 0.96 avg & \checkmark \\
O5 & HIGH violations escalated & 100\% & 100\% & \checkmark \\
O6 & Language-ID accuracy & $\geq$98\% & 98.2\% & \checkmark \\
O6 & Code-switching F1 & $\geq$92\% & 92.4\% & \checkmark \\
O7 & Question-count reduction & 40--50\% & --- & \textit{Pending final phase} \\
O8 & Median completion time & $\leq$17 min & --- & \textit{Pending final phase} \\
\bottomrule
\end{tabular}
}
\end{table}

% ------------------------------------------------------------
\section{Threats to Validity}
\label{sec:res-validity}
% ------------------------------------------------------------

\textbf{External validity.}
The pilot population (college students aged 18--25) is more digitally literate
and employed in simpler occupations than the general LFS target population.
Completion time and satisfaction advantages may attenuate with older workers
or informal-sector respondents.

\textbf{ISIC and ISCED coverage.}
The ISIC Rev.~4 classifier covers 100+ of 238 four-digit class entries; less
common industries may be classified less accurately.

\textbf{Statistical power for secondary outcomes.}
The $n=30$ pilot provides 80\% power for the primary outcome (completion time)
but is underpowered for smaller effects such as abandonment rate reduction
($n>100$ per arm needed).
